\documentclass[12pt]{article}
\usepackage[spanish]{babel}
\usepackage[utf8]{inputenc}
\usepackage{graphicx}
\usepackage{geometry}
\usepackage{parskip}
\usepackage{setspace}
\usepackage{titlesec}
\usepackage{tabularx}
\usepackage{booktabs}
\usepackage{hyperref}
\usepackage{float}
\usepackage{longtable}
\usepackage{array}
\usepackage{colortbl}
\usepackage{xcolor}

\geometry{a4paper, margin=2.5cm}
\setlength{\parindent}{0pt}
\renewcommand{\baselinestretch}{1.15}

% Configuración de títulos
\titleformat{\section}{\large\bfseries}{}{0em}{}
\titleformat{\subsection}{\normalsize\bfseries}{}{0em}{}
\titleformat{\subsubsection}{\normalsize\bfseries\itshape}{}{0em}{}

% Definición de colores (valores entre 0 y 1)
\definecolor{high}{rgb}{0.9,0.2,0.2}
\definecolor{medium}{rgb}{0.95,0.7,0.2}
\definecolor{low}{rgb}{0.2,0.7,0.3}

% Para referencias con sangría francesa
\usepackage{quoting}
\newenvironment{hangparas}{\begin{quoting}[hang=true, indent=0.5cm]}{\end{quoting}}

\begin{document}

% --- PORTADA ---
\begin{titlepage}
    \thispagestyle{empty}
    \begin{center}
        \textsc{\Large {\textbf{UNIVERSIDAD POLITÉCNICA DE AGUASCALIENTES}}}\\
        {\Large \uppercase{Ingeniería en Tecnologías de la Información e Innovación Digital} \\[1cm] }

        \begin{center}
            % APARTADO PARA IMAGEN 1: Logo UPA
            % \includegraphics[width=0.45\textwidth]{imagenes/LOGO UPA.png}
            % \hspace{0.5cm}
            % APARTADO PARA IMAGEN 2: Logo TIID
            % \includegraphics[width=0.35\textwidth]{imagenes/TIID UPA.png}\\
            [Logo UPA y TIID aquí]
        \end{center}

        \vspace{0.5cm}
        \rule{\linewidth}{0.3mm} \\
        { \Huge \bfseries Tarea 03: Aplicación de estándares\\[0.2cm]
          \large (ISO/IEC 25010, IEEE 730, CMMI, MoProSoft) }
        \rule{\linewidth}{0.3mm} \\

        \vspace{0.8cm}

        {\large \textbf{ASIGNATURA}}\\
        Estándares y Métricas para el Desarrollo de Software\\

        \vspace{0.4cm}

        {\large \textbf{DOCENTE}}\\
        Dr(c). Raul Alejandro Velasquez Ortiz \\

        \vspace{0.4cm}

        {\large \textbf{PROYECTO INTEGRADOR}}\\
        Outline - Wiki interna estilo Notion (Open Source maduro) \\

        \vspace{0.4cm}

        {\large \textbf{EQUIPO No. 4}}\\
        \vspace{0.4cm}
        \textbf{CHAIREZ ALBA OLIVIA LIZBETH - UP240594} \\
        \textbf{ESTRADA RAMÍREZ NOÉ EZEQUIEL - UP240770} \\
        \textbf{FACIO PALOS OMAR - UP240470} \\
        \textbf{JAUREGUI DE LA RIVA ALAN RICARDO - UP240960} \\
        \textbf{RUBIO VIRAMONTES BRIAN ODIN - UP250004}

        \vspace{0.4cm}

        {\large \textbf{GRUPO}}\\
        TIID05B

        \vfill
        %-- La repo si existe jsjsj--
        {\large \textbf{Repositorio Git:} \url{https://github.com/UP240594/outline-sqa-docs}}

        {\large Aguascalientes, México\\
        18 de mayo de 2026}
    \end{center}
\end{titlepage}

% --- ÍNDICE ---
\newpage
\tableofcontents
\newpage

% --- INTRODUCCIÓN ---
\section{Introducción}

\subsection{Contexto del proyecto}
El presente documento corresponde a la Tarea 03 de la asignatura Estándares y Métricas para el Desarrollo de Software, en la cual se aplican los conceptos teóricos de calidad de software al proyecto integrador del equipo: \textbf{Outline}, una aplicación web de tipo wiki interna con funcionalidades similares a Notion, de código abierto y con una comunidad activa.

Outline fue seleccionado del Catálogo de 30 alternativas (Semana 1) por su naturaleza madura (más de 5,000 estrellas en GitHub, licencia MIT) y por ser un caso realista para evaluar atributos de calidad como usabilidad, seguridad y fiabilidad. El proyecto será analizado desde la perspectiva de un equipo de desarrollo que busca adoptar, desplegar y posiblemente contribuir a este software en un entorno corporativo.


% --- Aun me faltan cosas, nose si este bien, pero eso se me ocurre, no puse las imagenes pq no se verian en git, igual subire el pdf ---

% -- Por cierto quedo empalmado pq sino no cabia  en la portada xd--
\end{document}